\documentclass[12pt]{article}
\usepackage[utf8]{inputenc}
\usepackage[T1]{fontenc}
\usepackage[margin=1in]{geometry}
\usepackage{setspace}
\setlength{\parindent}{2em}
\setlength{\parskip}{1.25em}
\renewcommand{\baselinestretch}{1.0}
\makeatletter
\newcommand{\@subtitle}{}
\newcommand{\subtitle}[1]{\renewcommand{\@subtitle}{#1}}
\renewcommand{\@maketitle}{%
  \newpage
  \null
  \vskip 1em%
  \begin{center}%
    {\LARGE \@title \par}%
    \ifx\@subtitle\@empty\else
      \vskip 0.5em%
      {\large \@subtitle \par}%
    \fi
  \end{center}%
  \par
  \vskip 1em}
\makeatother
\usepackage{xcolor}
\usepackage[colorlinks=true,urlcolor=blue!55!black,linkcolor=black]{hyperref}

\title{\textit{Corfield v. Coryell}, 6 F. Cas. 546 (C.C.E.D. Pa. 1823)}
\subtitle{Case Brief}
\author{}
\date{}

\begin{document}
\maketitle

\section{Memory Jogger}

In \href{https://www.courtlistener.com/opinion/9306338/corfield-v-coryell/}{\textit{Corfield v. Coryell}}, Justice Bushrod Washington, riding circuit, upheld New Jersey's ban on nonresidents harvesting oysters in state waters, famously defining Article IV's Privileges and Immunities Clause as protecting only ``fundamental'' rights belonging to citizens of all free governments---not a cotenancy in a state's common property.

\section{Detailed Case Facts}

In February 1819, plaintiff Corfield purchased a vessel named the \textit{Hiram} from one De Silver, holding a bill of sale and a federal coasting license on board (\href{https://opencasebook.org/casebooks/4501-constitutional-law/resources/12.6-corfield-v-coryell/}{6 F. Cas. at 547}). Corfield leased the \textit{Hiram} to a man named Hand, who in turn subleased the vessel to John Keene for \$10 per month. Sailing from Philadelphia with a crewmate named Courtney, Keene brought the \textit{Hiram} on May 15, 1821, to the oyster beds in Maurice River Cove---a shallow indentation of Delaware Bay along the southern New Jersey coast---where Keene and Courtney began raking and gathering oysters using a dredge.

On that day, a New Jersey vessel called the \textit{Independence}, carrying 40 to 50 men including Cumberland County magistrates, constables, and the collector of the port (armed with firearms and a small unmounted swivel), intercepted the \textit{Hiram} and ordered her to halt. Three or four men boarded the \textit{Hiram}, including defendant Coryell acting as prize master, and took the vessel up the Maurice River to Leesburg, New Jersey. Keene fled to avoid arrest by a local creditor, so process was served the next day on Courtney before two Cumberland County justices of the peace sitting at Leesburg. After examining witnesses and the vessel's papers, the magistrates condemned the \textit{Hiram} and her tackle and ordered her sold; the vessel was subsequently auctioned for \$10 (\href{https://opencasebook.org/casebooks/4501-constitutional-law/resources/12.6-corfield-v-coryell/}{6 F. Cas. at 547}).

Coryell justified the seizure under New Jersey's June 9, 1820 statute, ``An Act for the preservation of clams and oysters'' (\href{https://press-pubs.uchicago.edu/founders/documents/a4_2_1s18.html}{The Founders' Constitution}). Section 1 prohibited any person from raking or gathering oysters on New Jersey beds between May 1 and September 1 (\$10 penalty). Section 2 prohibited any resident or nonresident from dredging for oysters at any time (\$50 penalty). Section 6---the central provision at issue---made it unlawful for any person not ``an actual inhabitant and resident of this state'' to gather clams, oysters, or shells in any rivers, bays, or waters in New Jersey aboard any vessel not wholly owned by a New Jersey resident, imposing a \$10 penalty plus forfeiture of the vessel and equipment. Section 7 authorized any person to seize offending vessels and bring them before two county justices of the peace for condemnation and sale, splitting the net proceeds evenly between the county collector and the seizing prosecutors.

Historically, the dispute arose amid early-republic interstate friction over coastal fisheries and state territorial sovereignty in the Delaware Bay, just as the Marshall Court was defining the boundary between federal commercial authority and state police power in \textit{Gibbons v. Ogden}, 22 U.S. (9 Wheat.) 1 (1824).

\section{Procedural History}

Corfield filed a common-law action of trespass in the United States Circuit Court for the Eastern District of Pennsylvania against Coryell for illegally seizing, carrying away, and converting the \textit{Hiram}. Coryell pleaded not guilty, with leave to justify under the June 9, 1820 New Jersey oyster statute.

At the April Term 1823 jury trial, Justice Bushrod Washington (Associate Justice of the U.S. Supreme Court, riding circuit) noted the great importance of the constitutional and jurisdictional issues and recommended that the jury return a verdict for the plaintiff and assess damages, subject to the court's resolution of reserved questions of law. The jury returned a verdict for Corfield of \$560 subject to the reserved points (\href{https://opencasebook.org/casebooks/4501-constitutional-law/resources/12.6-corfield-v-coryell/}{6 F. Cas. at 549--50}). Counsel argued the reserved points of law at the October Term 1824, and the court took the case under advisement until the April Term 1825, when Justice Washington delivered the opinion and ordered judgment entered for the defendant Coryell. (Although the opinion was delivered in April 1825, the case is standardly cited as 1823 from the April Term 1823 trial caption in \textit{Federal Cases} and \textit{Washington's Circuit Court Reports}, 4 Wash. C.C. 371.) No appeal was taken to the Supreme Court.

\section{Judicial Votes}

\textbf{Single Opinion of the Circuit Court:} Associate Justice Bushrod Washington, riding circuit for the United States Circuit Court for the Eastern District of Pennsylvania, delivered the sole opinion of the court (Case No. 3,230, 6 F. Cas. 546, 4 Wash. C.C. 371).

\textbf{Concurrences:} None.

\textbf{Dissents:} None.

\section{Holding}

\textbf{Judgment:} Judgment entered for the defendant Coryell, setting aside the jury's provisional \$560 verdict for Corfield.

\textbf{Rule Announced:} (1) New Jersey's June 9, 1820 statute excluding nonresidents from harvesting oysters in state waters does not violate the Commerce Clause (Art. I, \S\ 8, cl. 3), the Privileges and Immunities Clause (Art. IV, \S\ 2, cl. 1), or the Admiralty and Maritime Jurisdiction Clause (Art. III, \S\ 2). (2) The Privileges and Immunities Clause protects only those rights that are ``in their nature, fundamental; which belong, of right, to the citizens of all free governments,'' and does not grant citizens of other states a cotenancy in the common property---such as oyster beds and fisheries---owned by a state for its own citizens. (3) Under common-law pleading rules, an owner who has leased a vessel to a bailee for an unexpired term lacks actual or constructive possession and therefore cannot maintain an action of trespass (only an action on the case).

\textbf{Holding vs. Dicta:} Structurally, \textit{Corfield} presents a classic holding-versus-dicta puzzle. Justice Washington explicitly declined to decide the third reserved question (whether Maurice River Cove lay within Cumberland County's statutory boundaries) because the fourth objection---that Corfield lacked possession to bring a trespass action while Keene's monthly lease was unexpired---was ``fatal'' (\href{https://press-pubs.uchicago.edu/founders/documents/a4_2_1s18.html}{6 F. Cas. at 555}). Because the pleading defect alone defeated the suit, strictly speaking every constitutional ruling in the opinion could be characterized as an alternative holding or extended dicta. Even treating the constitutional rulings as alternative holdings, Washington's celebrated catalogue of ``fundamental'' rights (life, liberty, property, travel, habeas corpus, court access, equal taxation, and the elective franchise) was manifestly dicta: the only Article IV point necessary to reject Corfield's claim was that harvesting oysters from state-owned beds is a non-fundamental common-property right. Yet that dicta became one of the most influential passages in American constitutional history, cited extensively during the framing of the Fourteenth Amendment and in \textit{The Slaughter-House Cases}, 83 U.S. (16 Wall.) 36 (1873) (\href{https://constitution.congress.gov/browse/essay/artIV-S2-C1-7/ALDE_00013783/}{Constitution Annotated}).

\section{Analysis of Opinions}

\subsection{Opinion of the Court (Circuit Justice Bushrod Washington)}

Justice Washington systematically addressed four sets of reserved objections (\href{https://press-pubs.uchicago.edu/founders/documents/a4_2_1s18.html}{6 F. Cas. at 550--55}):

\textbf{1. Commerce Clause (Art. I, \S\ 8, cl. 3):} Relying on Chief Justice Marshall's newly decided opinion in \textit{Gibbons v. Ogden}, 22 U.S. (9 Wheat.) 1 (1824), Washington defined interstate commerce as commercial intercourse and the navigation necessary to carry it on. While Congress holds exclusive authority over navigation and commercial intercourse on navigable waters (\textit{jus publicum}), that grant did not cede to the federal government a state's territory or its proprietary interest in the soil beneath its waters and the fisheries therein (\textit{jus privatum}). Citing Vattel's \textit{Law of Nations}, Washington held that a state retains full police and proprietary power to regulate the seasons, methods, and eligibility for harvesting its common-property oyster beds. Just as state inspection laws act on goods before they enter interstate commerce, New Jersey's statute did not obstruct free navigation or forbid buying and selling oysters once lawfully harvested; it merely regulated their removal from state-owned beds.

\textbf{2. Privileges and Immunities Clause (Art. IV, \S\ 2, cl. 1):} Turning to the Clause's command that ``the citizens of each state shall be entitled to all privileges and immunities of citizens in the several states,'' Washington rejected plaintiff's contention that nonresidents must be admitted to \textit{every} legal advantage enjoyed by local citizens. Instead, he confined the Clause to privileges and immunities ``which are, in their nature, fundamental; which belong, of right, to the citizens of all free governments; and which have, at all times, been enjoyed by the citizens of the several states which compose this Union.'' He grouped these under broad natural-rights headings---protection by government, enjoyment of life and liberty, acquiring and possessing property, and pursuing happiness and safety, subject to just restraints for the general good---and listed specific examples: passing through or residing in another state for trade, agriculture, or professional pursuits; claiming habeas corpus; suing in state courts; holding and disposing of real or personal property; exemption from higher taxes than state citizens pay; and the elective franchise as regulated by state law.

Washington then drew a sharp boundary between fundamental civil rights and state common property. Citing Grotius and Rutherfoth, he reasoned that where private fishery rights do not exist, the right to fish in state waters belongs to all citizens of the state as ``tenants in common.'' Construing Article IV to grant out-of-staters a cotenancy in New Jersey's oyster beds would risk exhausting and destroying a finite natural resource vital to the state's residents.

\textbf{3. Admiralty and Maritime Jurisdiction (Art. III, \S\ 2):} Plaintiff argued that taking oysters out of season with dredges was an ancient English admiralty misdemeanor, placing enforcement exclusively in federal admiralty courts. Washington rejected this claim on two grounds: first, the states never surrendered their exclusive pre-constitutional authority to regulate internal fisheries merely by granting admiralty jurisdiction to the federal judiciary; second, even if Congress could legislate over such offenses, Congress had passed no federal statute making oyster dredging in state waters a federal crime, leaving state enforcement intact.

\textbf{4. Territorial Boundaries, County Jurisdiction, and Form of Action:} Reviewing colonial royal grants from Charles II to the Duke of York (1664) and conveyances to Berkeley, Carteret, and William Penn, Washington concluded that although the original grants alone stopped at low-water mark, long-standing colonial usage, West Jersey statutes, and the 1783 Treaty of Peace confirmed New Jersey's title to Delaware Bay waters as against a private plaintiff asserting no conflicting title from Delaware. Washington expressed doubt as to whether Cumberland County's statutory boundary extended below low-water mark, but declined to decide that issue because the fourth objection was dispositive: because Corfield had leased the \textit{Hiram} to Hand and Keene for an unexpired monthly term, Corfield lacked both actual and constructive possession at the time of seizure and therefore could not maintain an action of trespass. (A fourth constitutional objection mentioned in argument---warrantless seizure under the Fourth Amendment---was not addressed in Washington's opinion.)

\section{Examples: Future Applications}

\begin{enumerate}
\item \textbf{Same-Side Application (Differential Fees for Recreational Big-Game Hunting Licenses):} Suppose Montana charges out-of-state recreational elk hunters \$225 for a hunting license while charging Montana residents \$9. Under \textit{Corfield}'s fundamental-rights limitation, the fee differential does not violate the Article IV Privileges and Immunities Clause because recreational sport hunting is not a fundamental right essential to the vitality of the Union or to earning a livelihood, mirroring the Supreme Court's holding in \textit{Baldwin v. Fish \& Game Commission of Montana}, 436 U.S. 371 (1978) (\href{https://constitution.congress.gov/browse/essay/artIV-S2-C1-7/ALDE_00013783/}{Constitution Annotated}).

\item \textbf{Same-Side Application (State Citizens-Only Access to Public Records Under a State FOIA Statute):} Suppose Virginia enacts a Freedom of Information Act granting Virginia citizens---but not nonresidents---a statutory right to request and inspect state government records. Under \textit{Corfield}, excluding nonresidents does not violate Article IV because a statutory public-records mechanism for holding state officials accountable to their electorate is not a fundamental right belonging to citizens of all free governments, as the Supreme Court held in \textit{McBurney v. Young}, 569 U.S. 221 (2013) (\href{https://constitution.congress.gov/browse/essay/artIV-S2-C1-7/ALDE_00013783/}{Constitution Annotated}).

\item \textbf{Opposite-Side Application (Prohibitive Licensing Fees on Nonresident Commercial Fishermen):} Suppose South Carolina charges nonresident commercial shrimp boats a \$2,500 licensing fee to fish for migratory shrimp in coastal waters, while charging resident boats only \$25. Unlike the stationary oyster beds in \textit{Corfield}, pursuing a commercial livelihood or common calling is squarely within Washington's list of fundamental privileges (``trade, agriculture, professional pursuits''), and a discriminatory fee designed to exclude nonresident commercial competitors violates Article IV, as held in \textit{Toomer v. Witsell}, 334 U.S. 385 (1948) (\href{https://constitution.congress.gov/browse/essay/artIV-S2-C1-7/ALDE_00013783/}{Constitution Annotated}).

\item \textbf{Opposite-Side Application (Residency Requirement for Admission to the State Bar):} Suppose New Hampshire refuses to admit an otherwise qualified Vermont resident who passed the New Hampshire bar exam solely because she does not reside in New Hampshire. Under \textit{Corfield}'s explicit protection of the right to pursue ``professional pursuits'' in another state, barring nonresidents from practicing law infringes a fundamental privilege without substantial justification and violates Article IV, as held in \textit{Supreme Court of New Hampshire v. Piper}, 470 U.S. 274 (1985) (\href{https://constitution.congress.gov/browse/essay/artIV-S2-C1-7/ALDE_00013783/}{Constitution Annotated}).

\item \textbf{Fence-Sitter Application (Municipal Hiring Preference for Local Residents on Publicly Funded Construction Projects):} Suppose a city mandates that at least 40\% of employees working on city-funded public construction contracts must be city residents. On one side, following \textit{Corfield}'s common-property rationale, the city can argue it is merely allocating its own public tax dollars and public-works contracts---analogous to common property---to its own taxpayers. On the other side, pursuing employment in the construction trades is a fundamental common calling under \textit{Corfield}, and excluding out-of-state workers burdens interstate commercial harmony. In \textit{United Building \& Construction Trades Council v. Mayor of Camden}, 465 U.S. 208 (1984), the Supreme Court held that such a preference implicates a fundamental privilege under Article IV, while remanding to evaluate whether the city had a ``substantial reason'' tied to local economic distress (\href{https://constitution.congress.gov/browse/essay/artIV-S2-C1-7/ALDE_00013783/}{Constitution Annotated}).
\end{enumerate}

\section{Critique}

Few circuit-court opinions have sparked more enduring constitutional debate than \textit{Corfield v. Coryell}, largely because Justice Washington's famous Article IV paragraph is notoriously ambiguous on a foundational question: does the Privileges and Immunities Clause guarantee \textit{substantive natural rights} against all states, or merely \textit{interstate nondiscrimination} with respect to fundamental rights?

\textbf{Originalist and Textualist Perspectives:} Originalist scholars are sharply divided over \textit{Corfield}'s legacy. Some originalists---including Randy Barnett, Richard Epstein, and John Eastman---read Washington's invocation of rights ``which are, in their nature, fundamental; which belong, of right, to the citizens of all free governments'' as confirming a natural-rights foundation that Reconstruction framers (such as Senator Jacob Howard) deliberately incorporated into the Fourteenth Amendment's Privileges or Immunities Clause, rendering the Supreme Court's narrow reading in \textit{The Slaughter-House Cases} (1873) unfaithful to original meaning (\href{https://scholarlycommons.law.northwestern.edu/cgi/viewcontent.cgi?article=1181&context=nulr}{Northwestern University Law Review}; \href{https://www.law.georgetown.edu/constitution-center/constitution/state-action/}{Georgetown Center for the Constitution}). Conversely, other originalists---such as Kurt Lash and Raoul Berger---argue that \textit{Corfield} has been widely misread: Article IV was an interstate comity and anti-discrimination clause (as later confirmed in \textit{Paul v. Virginia}, 75 U.S. (8 Wall.) 168 (1869)), not a roving license for federal judges to invent unenumerated substantive rights, and Representative John Bingham explicitly moved away from Article IV's language in drafting the final Fourteenth Amendment (\href{https://www.law.georgetown.edu/constitution-center/constitution/state-action/}{Georgetown Center for the Constitution}; \href{https://constitution.heritage.org/essays/a4-s2-c1/}{Heritage Guide to the Constitution}). Textualists also note a glaring internal anomaly in Washington's list: he included ``the elective franchise'' as a fundamental privilege, even though states have never been required to let out-of-state visitors vote in local elections, forcing him to qualify the item immediately with ``as regulated and established by the laws or constitution of the state.''

\textbf{Progressive and Doctrinal Critiques:} Progressive and modern doctrinal scholars criticize both \textit{Corfield}'s natural-law vagueness and its ``common property'' exception. First, Washington's appeal to rights belonging to ``all free governments'' offers no clear, objective boundary between fundamental and non-fundamental privileges, inviting unguided judicial policymaking---a weakness that eventually led the Supreme Court in \textit{Baldwin v. Fish \& Game Commission} (1978) to replace natural-law rhetoric with a functional test asking whether a right is ``basic to the maintenance or well-being of the Union.'' Second, modern scholars and courts have largely eroded \textit{Corfield}'s state-ownership (``cotenancy'') fiction for natural resources. In cases like \textit{Toomer v. Witsell} (1948) and \textit{Hicklin v. Orbeck}, 437 U.S. 518 (1978), the Court recognized that treating wildlife or natural resources as state ``common property'' can serve as a pretext for naked economic protectionism that undermines national economic unity (\href{https://constitution.congress.gov/browse/essay/artIV-S2-C1-7/ALDE_00013783/}{Constitution Annotated}). Finally, historians such as Gerard Magliocca have noted that \textit{Corfield}'s canonical status was largely a post-Civil War creation: before the 1860s, the opinion was rarely cited for its fundamental-rights passage, and its sweeping constitutional pronouncements were technically unnecessary once Washington held the trespass action barred by the bailment (\href{https://scholarship.law.nd.edu/cgi/viewcontent.cgi?article=4885&context=ndlr}{Notre Dame Law Review}; \href{https://harvardlawreview.org/print/vol-128/congresss-power-to-define-the-privileges-and-immunities-of-citizenship/}{Harvard Law Review}).

\textbf{Underlying Competing Values:} \textit{Corfield} crystallizes the permanent tension between national citizenship and state political community: a federal union requires citizens to move, trade, and own property freely across state lines, yet states cannot function as self-governing fiscal and ecological communities if they must share every locally funded benefit and finite public resource with the entire nation.

\section{Key Quotations}

\begin{quote}
The inquiry is, what are the privileges and immunities of citizens in the several states? We feel no hesitation in confining these expressions to those privileges and immunities which are, in their nature, fundamental; which belong, of right, to the citizens of all free governments; and which have, at all times, been enjoyed by the citizens of the several states which compose this Union, from the time of their becoming free, independent, and sovereign.
\end{quote}
--- Circuit Justice Bushrod Washington, \href{https://press-pubs.uchicago.edu/founders/documents/a4_2_1s18.html}{6 F. Cas. at 551}.

\begin{quote}
Protection by the government; the enjoyment of life and liberty, with the right to acquire and possess property of every kind, and to pursue and obtain happiness and safety; subject nevertheless to such restraints as the government may justly prescribe for the general good of the whole.
\end{quote}
--- Circuit Justice Bushrod Washington, \href{https://press-pubs.uchicago.edu/founders/documents/a4_2_1s18.html}{6 F. Cas. at 551--52}.

\begin{quote}
The right of a citizen of one state to pass through, or to reside in any other state, for purposes of trade, agriculture, professional pursuits, or otherwise; to claim the benefit of the writ of habeas corpus; to institute and maintain actions of any kind in the courts of the state; to take, hold and dispose of property, either real or personal; and an exemption from higher taxes or impositions than are paid by the other citizens of the state; may be mentioned as some of the particular privileges and immunities of citizens, which are clearly embraced by the general description of privileges deemed to be fundamental: to which may be added, the elective franchise, as regulated and established by the laws or constitution of the state in which it is to be exercised.
\end{quote}
--- Circuit Justice Bushrod Washington, \href{https://press-pubs.uchicago.edu/founders/documents/a4_2_1s18.html}{6 F. Cas. at 552}.

\begin{quote}
Where those private rights do not exist to the exclusion of the common right, that of fishing belongs to all the citizens or subjects of the state. It is the property of all; to be enjoyed by them in subordination to the laws which regulate its use. They may be considered as tenants in common of this property; and they are so exclusively entitled to the use of it, that it cannot be enjoyed by others without the tacit consent, or the express permission of the sovereign who has the power to regulate its use.
\end{quote}
--- Circuit Justice Bushrod Washington, \href{https://press-pubs.uchicago.edu/founders/documents/a4_2_1s18.html}{6 F. Cas. at 552}.

\begin{quote}
The grant to congress to regulate commerce on the navigable waters belonging to the several states, renders those waters the public property of the United States, for all the purposes of navigation and commercial intercourse; subject only to congressional regulation. But this grant contains no cession, either express or implied, of territory, or of public or private property.
\end{quote}
--- Circuit Justice Bushrod Washington, \href{https://press-pubs.uchicago.edu/founders/documents/a4_2_1s18.html}{6 F. Cas. at 551}.

\end{document}
